\documentclass[12pt]{article}
\usepackage[utf8]{inputenc}
\usepackage[T1]{fontenc}
\usepackage{hyperref}
\title{Jejara Naga Language Resources for K-12 Curriculum Development}
\author{Compiled by Claude Sonnet 5.5 (claude-sonnet-5-5) for Andr\'as Kornai}
\date{}
\begin{document}
\maketitle

\section{Introduction}
Jejara Naga, also called Para or Para Naga (ISO 639-3 pzn), is a Sino-Tibetan Naga language spoken in seven villages of Leshi Township, Hkamti District, Sagaing Region, Myanmar. Estimates range from about 1,460 to 1,800 speakers. It is used in most domains except official ones such as government, media and education. Its classification is disputed (Ao or Tangkhulic).

\section{Primary Educational Resources}
No school materials were found in the searches. A 2022 paper presents a Latin-based orthography for the previously unwritten language, which would be the basis for any primer.

\section{Linguistic and Reference Materials}
\begin{itemize}
\item Wikipedia, Para language: \url{https://en.wikipedia.org/wiki/Para_language}
\item Journal of the Southeast Asian Linguistics Society (Nov 2022), paper on the sociolinguistic background, phoneme inventory and Latin-based orthography of Para Naga: \url{https://doaj.org/article/d99795571ebd47359002269326853daf}
\item Endangered Languages Project entry for Jejara: \url{https://endangeredlanguages.com/lang/10415}
\end{itemize}
Search results also cite T. Barkman's 2014 Payap University MA thesis, \emph{A Descriptive Grammar of Jejara (Para Naga)}, but no URL for it was seen.

\section{Religious and Community Resources}
\begin{itemize}
\item Asia Harvest, Myanmar Para Naga profile: \url{https://www.operation.asiaharvest.org/MYANMAR-Para-Naga.pdf}
\item Joshua Project, Para Naga in Myanmar: \url{https://joshuaproject.net/people_groups/12659}
\end{itemize}
Bible portions are reported to have been published in 2022 (per search summary); no URL was seen.

\section{Digital and Audio Resources}
No audio, video or corpus resources were found.

\section{Curriculum Development Context}
The language is vigorously used by all community members but is not used in education. The new orthography, the grammar and Bible portions are the starting points for mother-tongue materials, to be developed with the community.

\end{document}